\section{Applications}
\subsection{Quadratic Sparsest Cut}

In the quadratic sparset cut one is required to find a subset of vertices $S$ that minimizes 
\[\frac{c\br{S, \bar{S}}}{\br{\spr{S}\cdot \spr{\bar{S}}}^2}.\]

Intuitively, this objective forces the solution to be more balanced. Observe, that if one assigns each pair of vertices $u,v$ a demand of $w\br{u,v}=1$, then we have that 
\[w\br{S, V}\cdot w\br{\bar{S}, V}=\br{\spr{S}+\spr{S}\cdot \spr{\bar{S}}}\cdot \br{\spr{\bar{S}}+\spr{S}\cdot \spr{\bar{S}}}=\spr{S}\cdot \spr{\bar{S}}+\spr{S}^2\cdot \spr{\bar{S}}+\spr{S}\cdot \spr{\bar{S}}^2+\br{\spr{S}\cdot \spr{\bar{S}}}^2.
\] 
Since $\spr{S},\spr{\bar{S}}\geq 1$, we get that 
\[\br{\spr{S}\cdot \spr{\bar{S}}}^2\leq w\br{S, V}\cdot w\br{\bar{S}, V}\leq 4\cdot\br{\spr{S}\cdot \spr{\bar{S}}}^2
\]
and the two objectives are equivalent up to a factor of $4$. Since $w$ is multiplicative, we obtain the following corollary:
\begin{corollary}
    There is a polynomial-time $O(\sqrt{\log n})$-approximation algorithm for the quadratic sparsest cut problem.
\end{corollary}
\subsection{Graph partitioning with demands}
Below, we show how to use the approximation algorithm for the minimum generalized conductance problem to obtain an approximation algorithm for the graph partitioning with demands. In this problem, we are given a graph $G=(V,E,c,w)$ and the goal is to find a subset of edges $C\subseteq E$ that minimizes $c\br{C}$ such that for every $H\in V-C$, we have $w\br{H}\leq w\br{V}/2$. We exhibit a bi-criteria approximation algorithm for this problem, which finds a subset $C$ such that $c\br{C}=O\br{\alpha}\cdot c\br{C^*}$ and for every $H\in V-C$, we have $w\br{H}\leq 4w\br{V}/5$, where $C^*$ is the optimal solution and $\alpha$ is the approximation ratio for the minimum generalized conductance problem. The analysis uses a generalization of the standard technique of dealing with balanced cut/graph partitioning problems using sparsest cut \cite{ExpanderFlowsGeometricEmbeddingsAndGraphPartitioning,ExpanderFlowsGeometricEmbeddingsAndGraphPartitioning}. However, due to the fact that we need to control both the demand in the components and the demand crossing the cut, a more careful analysis is required. The algorithm is simple:
\begin{enumerate}
    \item $C\gets\emptyset$
    \item $G'\gets \br{V, E,c,w}$
    \item While $w\br{V\br{G'}}> 4w\br{V}/5$ do
    \begin{itemize}
        \item Use the $\alpha$-approximation algorithm for the minimum generalized conductance problem to find a cut $(S,\bar{S})$ in $G'$, where $w\br{S}\leq w\br{\bar{S}}$.
        \item $C\gets C\cup E\br{S, \bar{S}}$
        \item Remove the vertices in $S$ from $G'$, i.e., $G'\gets G'\br{V\br{G'}\setminus S}$
    \end{itemize}
    \item Return $C$
\end{enumerate}

\begin{theorem}
    The above algorithm returns a subset of edges $C$ such that $c\br{C}=O\br{\alpha}\cdot c\br{C^*}$ and for every $H\in V-C$, we have $w\br{H}\leq 4w\br{V}/5$, where $C^*$ is the optimal solution and $\alpha$ is the approximation ratio for the minimum generalized conductance problem.
\end{theorem}
\begin{proof}
    Consider all components $H\in G-C$. If $H$ is $G'$ after the last iteration of the while loop, then by the stopping condition of this loop we have $w\br{H}\leq 4w\br{V}/5$. Otherwise, $H$ was the smaller side of the cut in some iteration of the while loop, and thus $w\br{H}\leq w\br{V}/2$. In either case, we have $w\br{H}\leq 4w\br{V}/5$.

    Now we want to reason about $c\br{C}$. Let $C^*$ be the optimal solution. We would like to partition components of $G-C^*$ into two vertex sets $A$ and $B$, such that $w\br{A, V}, w\br{B, V}\geq w\br{V}/4$. 
    \begin{lemma}
        For every graph $G$ and demand function $w$, there exists a partition of the components of $G-C^*$ into two vertex sets $A$ and $B$ such that $w\br{A, V}, w\br{B, V}\geq w\br{V}/4$.
    \end{lemma}
    \begin{proof}
    We have the following cases:
    \begin{enumerate}
        \item There exists $H\in G-C^*$ such that $w\br{H}\geq w\br{V}/4$. In this case, we can set $A=H$ and $B=V\setminus H$. We have 
        \[w\br{A, V}\geq w\br{A} \geq w\br{V}/4\]
         and 
        \[w\br{B, V}=w\br{V}-w\br{A}\geq w\br{V}/2\]
        since $w\br{A}=w\br{H}\leq w\br{V}/2$.
        \item For every $H\in G-C^*$, we have $w\br{H}\leq w\br{V}/4$. We proceed in iterations. Start with $A=\emptyset$ and keep adding vertices of components of $G-C^*$ to $A$ until $w\br{A}\geq w\br{V}/4$. Let $H$ be the last component added to $A$ and let $A'$ be $A$ one iteration before that. Consider the two cases:
        \begin{itemize}
            \item If $w\br{A', H}\geq w\br{G}/4$, then we set $A=A'$ and $B=V\setminus A'$. We have 
            \[
            w\br{A, V}\geq w\br{A, B} \geq w\br{A', H}\geq w\br{V}/4
            \]
            and similarily $w\br{B, V}\geq w\br{V}/4$. 
            \item Otherwise, $A$ stays unchanged and we set $B=V\setminus A$. We have $w\br{A, V}\geq w\br{V}/4$ and 
              \begin{align*}
              w\br{B, V}
              &= 
              w\br{G}-w\br{A}= w\br{G}-\br{w\br{A'}+w\br{A', H}+w\br{H}}
              \\&\geq 
              w\br{G}-3w\br{V}/4=w\br{V}/4.
            \end{align*}
        \end{itemize}
    \end{enumerate}
    This concludes the proof of the lemma.
  \end{proof}
    Therefore we can use the partition $\br{A,B}$ found using the above lemma to construct a lower bound on the cost of the solution found by our algorithm. Suppose that we are at some iteration when $w\br{V\br{G'}}> 4w\br{V}/5$. We have that 
    \[
    w\br{A\cap V\br{G'}, V\br{G'}}\geq w\br{A, V}-w\br{V\setminus V\br{G'}}\geq w\br{V}/4-w\br{V}/5=w\br{V}/20.
    \]
    Symmetrically, for $B$ we get that, $w\br{B\cap V\br{G'}, V\br{G'}}\geq w\br{V}/20$. Let $\br{S, \bar{S}}$ be the cut returned by the $\alpha$-approximation algorithm for the minimum generalized conductance problem in $G'$. We have
\begin{align*}
    \frac{c\br{S, \bar{S}}}{w\br{S, V\br{G'}}\cdot w\br{\bar{S}, V\br{G'}}} & \leq \alpha \cdot \frac{c\br{A, B}}{w\br{A\cap V\br{G'}, V\br{G'}}\cdot w\br{B\cap V\br{G'}, V\br{G'}}} \leq \frac{400\alpha\cdot c\br{A, B}}{w\br{V}^2}
\end{align*}

Now, by rearranging the above inequality we get:
\begin{align*}
    c\br{S, \bar{S}} & \leq \frac{400\alpha\cdot c\br{A, B}}{w\br{V}^2}\cdot w\br{S, V\br{G'}}\cdot w\br{\bar{S}, V\br{G'}} \leq  \frac{400\alpha\cdot c\br{C^*}\cdot w\br{S, V\br{G'}}}{w\br{V}}
\end{align*}
since trivially we have $c\br{A, B}\leq c\br{C^*}$ and $w\br{\bar{S}, V\br{G'}}\leq w\br{V\br{G'}}\leq w\br{V}$.

Let $\ell$ be the index of some iteration, let $G_{\ell}'$ be $G'$ at this iteration and let $\br{S_{\ell}, \bar{S_{\ell}}}$ be the cut at this iteration. By summing the above inequality over all iterations we get:
\begin{align*}
    c\br{C} & \leq \sum_{\ell} c\br{S_{\ell}, \bar{S_{\ell}}} \leq \frac{400\alpha\cdot c\br{C^*}}{w\br{V}}\cdot \sum_{\ell} w\br{S_{\ell}, V\br{G'_{\ell}}} \leq \frac{400\alpha\cdot c\br{C^*}}{w\br{V}}\cdot w\br{V} = 400\alpha\cdot c\br{C^*}
\end{align*}
where the last inequality follows since the sets $S_{\ell}$ are disjoint and thus $\sum_{\ell} w\br{S_{\ell}, V\br{G'_{\ell}}}\leq w\br{V}$.
\end{proof}

As an immediate corollary we get:
\begin{corollary}
    There exists an $O\br{\log n}$-approximation algorithm for the graph partitioning with demands on general graphs and an $O\br{1}$-approximation algorithm for the graph partitioning with demands on trees. Moreover, if the demand function is multiplicative, then there exists an $O\br{\sqrt{\log n}}$-approximation algorithm.
\end{corollary}

\subsection{Hierarchical clustering with demands}

In this problem one is given a graph $G=(V,E,c,w)$ and the goal is to find a hierarchical clustering tree of $G$. The clustering tree $T$ of $G$ is a binary tree where each node $u\in V\br{T}$ is associated with a subset of $S_u\subseteq V$ such that the subsets associated with children of $u$ form a partition of $S_u$. The root of $T$ is associated with $V$ and the leafs of $T$ are associated with singletons. The Dasgupta's clustering objective is to find a tree $T$ minimizing \[
\sum_{u\in V\br{T}} w\br{S_u}\cdot c\br{S_u, V\setminus S_u}.
\]
We show how to use the approximation algorithm for the graph partitioning with demands problem to obtain an approximation algorithm for the hierarchical clustering with demands problem. Note, that the assumption that the hierarchical clustering tree is binary can be relaxed, since it is easy to show that any hierarchical clustering tree can be transformed into a binary one without increasing the cost. The algorithm is simple: We start with $T$ being a single node associated with $V$. We keep applying the approximation algorithm for the graph partitioning with demands problem to partition the current cluster into smaller clusters and recurse on each of them until we get to singletons.

In order to show that the above algorithm achieves an $O\br{\alpha}$-approximation ratio, where $\alpha$ is the approximation ratio for the graph partitioning with demands problem, we use a generalized decomposition technique into levels due to \cite{ApproximateHierarchicalClusteringViaSparsestCutAndSpreadingMetrics}. Let $\cS_{t}$ be the set of all maximal clusters $S$ in $T^*$ such that $w\br{S}\leq t$. Observe that $\cS_t$ is a partition of $V$. Let $E_t^*$ be the set of all edges between clusters in $\cS_t$. For convenience let $E_0^*=E$. We call $E_t^*$ the cut at level $t$. We have the following lemma:
\begin{lemma}
    $c\br{T^*} = \sum_{t=1}^{w\br{V}-1}c\br{E_t^*}$
\end{lemma}
\begin{proof}
    Consider any edge $uv$. Let $S$ be the minimal cluster in $T^*$ containing both $u$ and $v$. We get that the contribution o $uv$ to $c\br{T^*}$ is $w\br{S}\cdot c\br{uv}$. Moreover, we have that $uv\in E_t^*$ for every $t$ such that $t < w\br{S}$, so it contributes $c\br{uv}$ to $c\br{E_t^*}$ for every $t < w\br{S}$ which gives the claim. 
\end{proof}

We also use the following upper bound on the $c\br{T^*}$ which is a direct consequence of the previous lemma:
\begin{corollary}
    $2\cdot c\br{T^*} = 2\cdot \sum_{t=1}^{w\br{V}-1}c\br{E_t^*}\geq \sum_{t=0}^{w\br{V}}c\br{E_{\fl{t/2}}^*}$
\end{corollary}

Fix some non-leaf cluster $H$ in the tree $T$ returned by the algorithm and let $C_H$ be the cut returned by the graph partitioning with demands algorithm. Let $r_H=w\br{H}$ and $s_H$ be the size of the largest child of $H$. We have $s_H\leq 4r_H/5$. Edges cut in $H$ contribute $r_H\cdot c\br{C_H}$ to the objective function. We wish to charge this cost to the edges cut in $E_{r_H/2}^*$, restricted to $H$. We have the following bound:
\begin{align*}
    \br{r_H-s_H}\cdot c\br{E_{r_H/2}^*\cap H}\leq \sum_{t=s_H+1}^{r_H}c\br{E_{\fl{t/2}}^*\cap H}
\end{align*}
since as the level $t$ decreases, the more edges belong to $E_t^*$. 

We get the contribution of $C_H$ to the objective function is at most:
\begin{align*}
    r_H\cdot c\br{C_H} \br{r_H-s_H}\cdot c\br{E_{r_H/2}^*\cap H}\leq 5\cdot \sum_{t=s_H+1}^{r_H}c\br{E_{\fl{t/2}}^*\cap H}
\end{align*}
where the first inequality is due to the fact that $r_H-s_H\geq r_H/5$.

Now, since all of child clusters of $H$ have size at most $s_H$, they do not contribute to $c\br{E_{t}^*}$ for $t > s_H$. Combining this observation with the fact that the clusters in $T$ contributing to a given level $t$ form a partition of $V$, and thus are disjoint we get that:
\begin{align*}
    \sum_{H} \sum_{t=s_H+1}^{r_H}c\br{E_{t}^*\cap H} \leq \sum_{t=0}^{w\br{V}} c\br{E_{\fl{t/2}}^*}
\end{align*}

Thus by summing up over the conitribution of all cuts $C_H$ we get that the cost of the tree returned by the algorithm is at most:
\begin{align*}
    \sum_{H} r_H\cdot c\br{C_H} & \leq 5\cdot \sum_{H} \sum_{t=s_H+1}^{r_H}c\br{E_{t}^*\cap H} \leq 5\cdot \sum_{t=0}^{w\br{V}} c\br{E_{\fl{t/2}}^*}\leq 10\cdot c\br{T^*}
\end{align*}
where the last inequality is by the previous corollary.

We get the following theorem:
\begin{theorem}
    The above algorithm achieves an $O\br{\alpha}$-approximation ratio, where $\alpha$ is the approximation ratio for the graph partitioning with demands problem.
\end{theorem}

And as an immediate corollary we get:
\begin{corollary}
    There exists an $O\br{\log n}$-approximation algorithm for the hierarchical clustering with demands problem on general graphs and an $O\br{1}$-approximation algorithm for the hierarchical clustering with demands problem on trees. Moreover, if the demand function is multiplicative, then there exists an $O\br{\sqrt{\log n}}$-approximation algorithm.
\end{corollary}

